%% Correction audit de veracite - 2026-09-20 : citations de rapport de jury sans
%% support textuel dans le corpus (rapports 2018-2025) remplacees par des verbatims
%% verifies ou supprimees. Detail ci-dessous.
%% l.269 (ex) : citation "jury 2021" sur la lecture des courbes / correction en
%%   temperature : INTROUVABLE dans tout le corpus -> remplacee par RJ2021 l.759 (verbatim).
%% l.295 (ex) : citation "jury 2021" sur la justification du choix de l'onduleur :
%%   INTROUVABLE -> remplacee par RJ2021 l.757 (verbatim).
\expandafter\ifx\csname CAPETINCLUDEMODE\endcsname\relax
\documentclass[11pt,a4paper]{article}
% preamble.tex — packages et commandes communes a toutes les fiches CAPET SII IE
\input{glyphtounicode}
\pdfgentounicode=1
\usepackage{cmap}
\usepackage[T1]{fontenc}
\usepackage[utf8]{inputenc}
\usepackage[french]{babel}
% babel[french] active a \begin{document} des raccourcis de ponctuation (:;!?)
% qui cassent la syntaxe to[...] de circuitikz (erreur "+ or - expected") ;
% on les desactive, ce qui ne change que l'espacement fin avant ces signes.
\AtBeginDocument{\shorthandoff{:;!?}}
\usepackage[a4paper, top=1.6cm, bottom=1.6cm, left=1.4cm, right=1.4cm]{geometry}
\setlength{\headheight}{14.5pt}
\usepackage{amsmath, amssymb, amsfonts}
\usepackage{siunitx}
\sisetup{locale=FR, output-decimal-marker={,}, inter-unit-product=\ensuremath{{}\cdot{}},
         group-minimum-digits=4}
% Unites non fournies par siunitx, declarees ici une fois pour toutes (fiches + exercices).
% Les gardes \ifdefined evitent tout conflit avec une redeclaration locale restee dans un
% fichier : ne pas retirer les gardes, ne pas redeclarer ces unites dans les fichiers.
\DeclareSIUnit{\jour}{j}
\ifdefined\var\else\DeclareSIUnit{\var}{var}\fi   % puissance reactive
\ifdefined\vad\else\DeclareSIUnit{\vad}{VAD}\fi   % puissance deformante
\ifdefined\tr\else\DeclareSIUnit{\tr}{tr}\fi      % tour (tr/min)
\usepackage{array, booktabs, multirow, tabularx}
\usepackage{xcolor}
\usepackage[most, breakable]{tcolorbox}
\usepackage{multicol}
\usepackage{enumitem}
\usepackage{mathtools}
\usepackage{lmodern}
% Refonte visuelle 2026 (direction "moderne coloré") : Lato pour le texte,
% maths sans-serif assorties via newtxsf (charge apres amsmath/amssymb).
\usepackage[default]{lato}
\usepackage{newtxsf}
\usepackage{microtype}
\usepackage{fancyhdr}
\usepackage{tikz}
\usetikzlibrary{positioning,arrows.meta}
\usepackage{circuitikz}
\usepackage{lastpage}

% ── Couleurs par domaine (palette modernisée 2026) ────────────────────
\definecolor{colElec}{RGB}{29, 78, 216}        % Électrotechnique — bleu vif
\definecolor{colElecLight}{RGB}{219, 234, 254} % fond bleu clair
\definecolor{colAuto}{RGB}{21, 128, 61}        % Automatique — vert
\definecolor{colMeca}{RGB}{194, 65, 12}        % Mécanique — orange
\definecolor{colTherm}{RGB}{185, 28, 28}       % Thermique — rouge
\definecolor{colQT}{RGB}{232, 244, 233}        % (hérité, plus utilisé par questiontype)
\definecolor{colJury}{RGB}{255, 248, 225}      % fond jury (crème)
\definecolor{colJuryFrame}{RGB}{202, 138, 4}   % bordure/bandeau jury (ambre)
\definecolor{colGrayLight}{RGB}{245, 245, 247} % fond gris neutre
\definecolor{colSeparator}{RGB}{180, 180, 180} % gris séparateur
\definecolor{colExo}{RGB}{8, 145, 178}          % Banque d'exercices — cyan (distinct des fiches)
% ── Encre et gris de la refonte ───────────────────────────────────────
\definecolor{colInk}{RGB}{23, 32, 42}          % texte courant (quasi-noir)
\definecolor{colSlate}{RGB}{83, 98, 118}       % gris ardoise (méta, en-têtes)
\definecolor{colRule}{RGB}{203, 213, 225}      % filets clairs
\definecolor{colJuryBg}{RGB}{255, 248, 225}    % fond crème jury
\definecolor{colPresence}{RGB}{83, 98, 118}    % bloc 5 (présence dans les sujets)
\AtBeginDocument{\color{colInk}}

% ── Style encadré exercice + corrigé (banque d'exercices) ─────────────
\tcbset{
  exoBox/.style={
    breakable,
    enhanced,
    colback=colExo!8!white,
    colframe=colExo,
    colbacktitle=colExo,
    coltitle=white,
    fonttitle=\bfseries\small,
    left=4pt, right=4pt, top=3pt, bottom=3pt,
    before skip=5pt, after skip=5pt,
  }
}

% ── Bascule enonce / corrige (banque d'exercices) ─────────────────────
% \CAPETCORRIGE est defini en ligne de commande (pdflatex "\def\CAPETCORRIGE{1}\input{...}")
% pour compiler la variante corriges d'un meme fichier source ; par defaut (non defini),
% la compilation produit la variante enonces. \CAPETSHOW{enonce}{corrige} choisit la branche.
%
% TROIS MODES (refonte 2026-09-20, manuel de revision) :
%   - aucun drapeau   -> imprime l'enonce            (variante enonces, INCHANGE)
%   - \CAPETCORRIGE   -> imprime le corrige          (variante corriges, INCHANGE)
%   - \CAPETMANUEL    -> imprime l'enonce ET met le corrige de cote pour l'annexe
%
% Le mode manuel stocke chaque corrige dans une macro (\csgdef), sans ecriture
% fichier ni re-tokenisation : le contenu (circuitikz, accents, %) traverse intact.
% L'annexe les rejoue ensuite dans l'ordre via \CAPETANNEXE.
\usepackage{etoolbox}

\newcounter{CAPETexo}          % numero continu d'exercice dans le manuel
\setcounter{CAPETexo}{0}
\newcounter{CAPETannexe}       % compteur de relecture en annexe

% Titre courant de la notion en cours (pose par \CAPETNOTION), sert d'intertitre annexe
\newcommand{\CAPETnotioncode}{}
\newcommand{\CAPETnotiontitre}{}

% \CAPETNOTION{code}{titre} : ouvre une notion dans le manuel (memorise le libelle,
% et marque le rang du premier exercice de la notion pour l'intertitre d'annexe).
\newcommand{\CAPETNOTION}[2]{%
  \renewcommand{\CAPETnotioncode}{#1}%
  \renewcommand{\CAPETnotiontitre}{#2}%
  \csgdef{CAPETnotionstart\the\numexpr\value{CAPETexo}+1\relax}{#1 -- #2}%
}

% ── \degree robuste ───────────────────────────────────────────────────
% Certains fichiers d'exercices utilisent \degree. Ni le preambule ni siunitx
% ne le definissent, et \providecommand local ne suffit pas : sous hyperref,
% \degree existe deja comme commande factice (\ERROR), donc \providecommand
% s'abstient et l'appel echoue a la compilation du manuel (erreur invisible
% en compilation isolee). On le (re)definit donc inconditionnellement ici,
% apres le chargement des packages, pour tout le corpus.
\AtBeginDocument{\providecommand{\degree}{}\renewcommand{\degree}{\ensuremath{^\circ}}}

% Libelle decoratif : les fichiers d'exercices utilisent aussi \CAPETSHOW dans
% leur \fancyhead et leur bloc de titre, pour afficher "(enonces)" ou "(corriges)".
% Ces appels-la ne sont PAS des exercices : en mode manuel ils ne doivent ni
% consommer un numero, ni alimenter l'annexe. On les reconnait a leur premier
% argument, strictement egal au libelle ci-dessous.
\newcommand{\CAPETLIBELLEENONCES}{(\'enonc\'es)}

\ifdefined\CAPETMANUEL
  % --- Mode manuel : enonce en place, corrige differe vers l'annexe ---
  \newcommand{\CAPETSHOW}[2]{%
    \def\CAPETargun{#1}%
    \ifx\CAPETargun\CAPETLIBELLEENONCES
      #1%            appel decoratif : rendu tel quel, hors numerotation
    \else
      \CAPETSHOWEXO{#1}{#2}%
    \fi
  }
  \newcommand{\CAPETSHOWEXO}[2]{%
    \stepcounter{CAPETexo}%
    \csgdef{CAPETcorrige\theCAPETexo}{#2}%
    \edef\CAPETtmpnum{\theCAPETexo}%
    #1%
    \par\vspace{2pt}%
    {\footnotesize\itshape\color{colExo}%
     \hyperref[capet:cor:\CAPETtmpnum]{Corrige de l'exercice \CAPETtmpnum\ $\rightarrow$ p.~\pageref{capet:cor:\CAPETtmpnum}}%
     \label{capet:eno:\CAPETtmpnum}}%
  }
  % \CAPETANNEXE : rejoue tous les corriges captures, dans l'ordre, avec intertitres
  \newcommand{\CAPETANNEXE}{%
    \setcounter{CAPETannexe}{0}%
    \loop\ifnum\value{CAPETannexe}<\value{CAPETexo}%
      \stepcounter{CAPETannexe}%
      \edef\CAPETcurnum{\theCAPETannexe}%
      \ifcsname CAPETnotionstart\CAPETcurnum\endcsname
        \par\medskip
        {\normalsize\bfseries\color{colExo}\csuse{CAPETnotionstart\CAPETcurnum}}%
        \par\nobreak\smallskip
      \fi
      \begin{tcolorbox}[exoBox, breakable,
        title={Corrige de l'exercice \CAPETcurnum\
               \hfill\normalfont\footnotesize
               \hyperref[capet:eno:\CAPETcurnum]{enonce p.~\pageref{capet:eno:\CAPETcurnum}}}]%
      \label{capet:cor:\CAPETcurnum}%
      \csuse{CAPETcorrige\CAPETcurnum}%
      \end{tcolorbox}%
    \repeat
  }
\else
  % --- Modes historiques, comportement strictement inchange ---
  \newcommand{\CAPETSHOW}[2]{\ifdefined\CAPETCORRIGE #2\else #1\fi}
  \newcommand{\CAPETANNEXE}{}
\fi

% ── Difficulté 1 à 3 (étoiles pleines, amssymb) ───────────────────────
\newcommand{\difficulte}[1]{%
  \ifcase#1\relax
  \or $\bigstar$\;niveau 1/3%
  \or $\bigstar\bigstar$\;niveau 2/3%
  \or $\bigstar\bigstar\bigstar$\;niveau 3/3%
  \fi
}

% ── Style encadré extrait de document technique (banque d'exercices) ──
\tcbset{
  docTech/.style={
    breakable,
    enhanced,
    colback=colGrayLight,
    colframe=colSeparator,
    fonttitle=\bfseries\small\itshape,
    coltitle=black,
    colbacktitle=colGrayLight,
    left=4pt, right=4pt, top=3pt, bottom=3pt,
    before skip=4pt, after skip=4pt,
  }
}

% ── Style tcolorbox de section principale ────────────────────────────
% Refonte 2026 : bandeau titre coloré arrondi + fond très pâle, sans cadre.
\tcbset{
  ficheSection/.style={
    breakable, enhanced jigsaw,
    frame hidden, boxrule=0pt, arc=1.8mm,
    colback=#1!4!white,
    colbacktitle=#1,
    coltitle=white,
    fonttitle=\sffamily\bfseries,
    left=7pt, right=7pt, top=5pt, bottom=5pt,
    toptitle=3pt, bottomtitle=3pt,
    before skip=10pt, after skip=8pt,
  }
}

% ── Style encadré question type ──────────────────────────────────────
% Carte blanche, bandeau sombre, liseré très clair.
\tcbset{
  questiontype/.style={
    breakable, enhanced jigsaw,
    colback=white,
    frame hidden, boxrule=0pt, arc=1.8mm,
    colbacktitle=colElec!78!black,
    coltitle=white,
    fonttitle=\sffamily\bfseries\small,
    left=7pt, right=7pt, top=4pt, bottom=4pt,
    toptitle=2.5pt, bottomtitle=2.5pt,
    before skip=8pt, after skip=8pt,
    borderline={0.6pt}{0pt}{colElec!12!white},
  }
}

% ── Style encadré jury ───────────────────────────────────────────────
% Bandeau ambre, fond crème, sans cadre.
\tcbset{
  jury/.style={
    breakable, enhanced jigsaw,
    colback=colJuryBg,
    frame hidden, boxrule=0pt, arc=1.8mm,
    colbacktitle=colJuryFrame,
    coltitle=white,
    fonttitle=\sffamily\bfseries,
    left=7pt, right=7pt, top=4pt, bottom=4pt,
    toptitle=3pt, bottomtitle=3pt,
    before skip=10pt, after skip=8pt,
  }
}

% ── En-tête / pied de page communs (refonte 2026) ─────────────────────
% #1 code fiche, #2 titre court, #3 couleur domaine
\newcommand{\ficheHeaders}[3]{%
  \fancyhead[L]{\small\sffamily{\bfseries\textcolor{#3}{#1}}\hspace{0.7em}\textcolor{colSlate}{#2}}%
  \fancyhead[R]{\small\sffamily\textcolor{colSlate}{CAPET SII IE\, --- \,Session 2026}}%
  \fancyfoot[C]{\small\sffamily\textcolor{colSlate}{\thepage\,/\,\pageref{LastPage}}}%
  \fancyhead[C]{}\fancyfoot[L]{}\fancyfoot[R]{}%
}

% ── Bloc en-tête de fiche (BLOC 0, refonte 2026) ──────────────────────
% #1 code  #2 titre  #3 couleur domaine  #4 domaine  #5 ligne présence  #6 sous-thèmes
\newcommand{\ficheEntete}[6]{%
  \begin{tcolorbox}[enhanced, breakable,
    frame hidden, boxrule=0pt, arc=2.5mm,
    interior style={left color=#3, right color=#3!75!black},
    left=9pt, right=9pt, top=7pt, bottom=7pt,
    before skip=0pt, after skip=6pt]
    {\color{white}\sffamily\bfseries\LARGE #1\quad #2}\\[4pt]
    {\color{#3!22!white}\sffamily\small \MakeUppercase{#4}\quad\textbar\quad #5}\\[1pt]
    {\color{#3!35!white}\sffamily\footnotesize Sous-thèmes : #6}
  \end{tcolorbox}%
}

% ── Mise en page ─────────────────────────────────────────────────────
\pagestyle{fancy}
\fancyhf{}
\renewcommand{\headrulewidth}{0.5pt}
\renewcommand{\footrulewidth}{0pt}
\AtBeginDocument{%
  \renewcommand{\headrule}{{\color{colRule}\hrule width\headwidth height 0.5pt}}%
}

\setlist[itemize]{leftmargin=1.2em, itemsep=1pt, topsep=2pt, parsep=0pt}
\setlist[enumerate]{leftmargin=1.2em, itemsep=3pt, topsep=2pt, parsep=0pt}

\setlength{\parskip}{2pt}
\setlength{\columnsep}{10pt}

\def\CAPETBEGINFILE{\begin{document}}
\def\CAPETENDFILE{\end{document}}
\else
\def\CAPETBEGINFILE{}
\def\CAPETENDFILE{}
\fi

\ficheHeaders{T16}{Énergies renouvelables}{colElec}

\CAPETBEGINFILE

% -----------------------------------------------------------------------
%  BLOC EN-TETE
% -----------------------------------------------------------------------

\ficheEntete{T16}{Énergies renouvelables : photovoltaïque, hydraulique, éolien}{colElec}%
  {électrotechnique - Systèmes énergétiques}%
  {\textbf{Présence :} \textbf{6/13} sujets (2013-2026) - thème \emph{fréquent}}%
  {modèle panneau PV ($I_{sc}$, $V_{oc}$, MPP) - nombre de panneaux et superficie - MPPT algorithme P\&O - hacheur abaisseur - micro-hydraulique (STEP) - éolien (loi de Betz, DFIG)}

% -----------------------------------------------------------------------
%  1. RAPPELS DE COURS
% -----------------------------------------------------------------------

\begin{tcolorbox}[
  ficheSection=colElec,
  title={1. Rappels de cours},
]

% -- 1.1 Panneau photovoltaeque ----------------------------------------
\textbf{1.1 Modèle électrique d'un panneau PV (STC : $1000\,\mathrm{W\,m^{-2}}$, $25^\circ\mathrm{C}$)}

\smallskip
\begin{center}
\renewcommand{\arraystretch}{1.3}
\begin{tabular}{lllll}
\toprule
\textbf{Paramètre} & \textbf{Notation} & \textbf{Valeur typ.} & \textbf{En série} & \textbf{En parallèle} \\
\midrule
Puissance crête & $P_{crete}$ & 320\,W & - & - \\
Courant court-circuit & $I_{sc}$ & 9{,}76\,A & $I_{sc}$ inchangé & additionné \\
Tension circuit ouvert & $V_{oc}$ & 40{,}96\,V & additionnée & $V_{oc}$ inchangé \\
Courant MPP & $I_{MPP}$ & 9{,}34\,A & idem & additionné \\
Tension MPP & $V_{MPP}$ & 34{,}32\,V & additionnée & idem \\
\bottomrule
\end{tabular}
\end{center}
\textbf{Correction en température :} $+1^\circ\mathrm{C} \to -0{,}4\,\%$ de $P_{crete}$ (Si-mono). À $50^\circ\mathrm{C}$ (STC + $25^\circ\mathrm{C}$) : $-10\,\%$ de puissance.

Puissance du champ : $P_{champ} = N_{panneaux} \times P_{crete} \times G/G_0$ avec $G_0 = 1000\,\mathrm{W\!\cdot\!m^{-2}}$.

\medskip\hrule\vspace{4pt}

% -- 1.2 MPPT ---------------------------------------------------------
\textbf{1.2 MPPT - Maximum Power Point Tracking}

\smallskip
Commande ajustant le rapport cyclique $\alpha$ du hacheur abaisseur pour maintenir le panneau au point de puissance maximale ($V_{MPP}$, $I_{MPP}$).

\smallskip
\begin{center}
\begin{tikzpicture}[scale=0.85, font=\footnotesize]
  \draw[->] (0,0) -- (6.3,0) node[right] {$U$};
  \draw[->] (0,-0.3) -- (0,2.6) node[above] {$P$};
  \draw[thick, colElec, domain=0.15:5.9, samples=60, smooth] plot (\x, {2.2*sin((\x/6)*180)});
  \draw[dashed] (3.2,0) -- (3.2,2.18);
  \node[below] at (3.2,0) {\scriptsize $U_{MPP}$};
  \node[above] at (3.2,2.2) {\scriptsize MPP};
  \draw[dashed] (1.3,0) -- (1.3,1.32);
  \node[below] at (1.3,0) {\scriptsize $U_1$};
  \draw[dashed] (4.6,0) -- (4.6,1.7);
  \node[below] at (4.6,0) {\scriptsize $U_2$};
\end{tikzpicture}
\end{center}

\textbf{Algorithme perturbation et observation (P\&O) :}
\begin{itemize}[itemsep=1pt]
  \item Imposer $U + \Delta U$ au panneau
  \item Si $\Delta P > 0$ : on se rapproche du MPP - continuer dans le même sens
  \item Si $\Delta P < 0$ : on s'éloigne - inverser le sens de perturbation
  \item Répéter jusqu'à $\Delta P \approx 0$ (convergence vers le MPP)
\end{itemize}

Rapport cyclique hacheur abaisseur : $\alpha = E_b / V_{PV}$ (la tension PV, élevée, est abaissée vers la tension du bus continu $E_b$ ; on a bien $0 < \alpha < 1$).

\medskip\hrule\vspace{4pt}

% -- 1.3 Connexion sur onduleur ----------------------------------------
\textbf{1.3 Dimensionnement de la connexion sur onduleur}

\smallskip
\[
  N_{serie/entree} = \frac{V_{MPP\,max\,onduleur}}{V_{MPP\,panneau}}
  \qquad;\qquad
  N_{onduleurs} = \frac{N_{panneaux\,total}}{N_{entrees/onduleur} \times N_{serie/entree}}
\]
Vérifier que $V_{MPP}$ et $I_{MPP}$ calculés sont dans les plages min-max de l'onduleur.

\medskip\hrule\vspace{4pt}

% -- 1.4 eolien -------------------------------------------------------
\textbf{1.4 Éolien - loi de Betz}

\smallskip
\[
  P_{extrait} = C_p \cdot \frac{1}{2}\,\rho_{air}\,A_{rotor}\,v^3
  \qquad;\qquad
  C_p^{max} = \frac{16}{27} \approx 0{,}593 \quad \text{(limite de Betz)}
\]
$A_{rotor} = \pi R^2$, $v$ = vitesse du vent. Si $v$ double, $P$ est multipliée par \textbf{8}.

\end{tcolorbox}

% -----------------------------------------------------------------------
%  2. FORMULAIRE SYNTHETIQUE
% -----------------------------------------------------------------------

\begin{tcolorbox}[
  ficheSection=colElec,
  title={2. Formulaire synthétique},
]
\small
\begin{multicols}{2}
\renewcommand{\arraystretch}{1.45}
\begin{tabular}{@{}p{4.5cm}ll@{}}
\toprule
\textbf{Formule} & \textbf{Usage} & \textbf{U.} \\
\midrule
\multicolumn{3}{@{}l}{\footnotesize\textit{Champ PV}} \\
$P_{champ} = N \cdot P_{crete} \cdot G/G_0$ & Puissance & W \\
$N = P_{souhaitee}/P_{crete}$ & Nb panneaux & - \\
$V_{MPP,chaine} = N_s V_{MPP}$ & Tension série & V \\
\midrule
\multicolumn{3}{@{}l}{\footnotesize\textit{MPPT}} \\
$\alpha = E_b/V_{PV}$ (abaisseur) & Rapp. cyclique & - \\
P\&O : si $\Delta P>0$ - même sens & - & - \\
P\&O : si $\Delta P<0$ - inverser & - & - \\
\bottomrule
\end{tabular}

\columnbreak

\renewcommand{\arraystretch}{1.45}
\begin{tabular}{@{}p{4.5cm}ll@{}}
\toprule
\textbf{Formule} & \textbf{Usage} & \textbf{U.} \\
\midrule
\multicolumn{3}{@{}l}{\footnotesize\textit{Hydraulique}} \\
$E_p = \rho V g h$ & Énergie pot. & J \\
$E_{elec} = E_p\cdot\eta_{turb}$ & Turbinage & W\,h \\
$T = V/Q$ & Durée pomp. & h \\
\midrule
\multicolumn{3}{@{}l}{\footnotesize\textit{Éolien}} \\
$P = C_p \cdot \frac{1}{2}\rho A v^3$ & Puissance & W \\
$C_p^{max} = 16/27$ & Betz & - \\
\midrule
\multicolumn{3}{@{}l}{\footnotesize\textit{Données 2021}} \\
$N = 782$, $S = 1300\,\mathrm{m^2}$ & PV 250\,kW & - \\
$N_s = 26$ panneaux/entrée & Onduleur & - \\
\bottomrule
\end{tabular}
\end{multicols}
\normalsize
\end{tcolorbox}

% -----------------------------------------------------------------------
%  3. QUESTIONS TYPES
% -----------------------------------------------------------------------

\begin{tcolorbox}[
  ficheSection=colElec,
  title={3. Questions types - extraites des sujets réels},
]

% -- QT1 -------------------------------------------------------------

\begin{tcolorbox}[questiontype,
  title={QT1 - Nombre de panneaux et paramétrage cellule \hfill\normalfont\footnotesize E2 disc. 2021 (Q11-Q12)}
]
\textbf{Énoncé.}
Panneau DMH320M6 : $P_{crete} = 320\,\mathrm{W}$, 60 cellules en série, $I_{sc} = 9{,}76\,\mathrm{A}$, $V_{oc} = 40{,}96\,\mathrm{V}$, surface 1{,}66\,m². Puissance crête souhaitée : 250\,kW.
(a) Calculer le nombre de panneaux et la superficie. (b) Indiquer $I_{sc,cell}$ et $V_{oc,cell}$ pour une cellule.

\tcbline

\textbf{Corrigé.}

\textbf{Étape 1 (a).} $N = 250\,000 / 320 = 782\,\text{panneaux}$ ; superficie $= 782 \times 1{,}66 = \mathbf{1298\,\mathrm{m^2}} \approx 1300\,\mathrm{m^2}$.

\textbf{Étape 2 (b).} Cellules en série - courant identique - $I_{sc,cell} = I_{sc,panneau} = \mathbf{9{,}76\,\mathrm{A}}$.\\
Tensions s'ajoutent - $V_{oc,cell} = V_{oc,panneau}/60 = 40{,}96/60 = \mathbf{0{,}683\,\mathrm{V}}$.

\textbf{Conclusion.} Une surface d'environ 1300\,m² est nécessaire pour atteindre les 250\,kW crête visés : c'est cette surface, comparée à celle disponible sur site, qui valide (ou non) la faisabilité de l'installation.
\end{tcolorbox}

\vspace{4pt}

% -- QT2 -------------------------------------------------------------

\begin{tcolorbox}[questiontype,
  title={QT2 - Dimensionnement entrées onduleur et vérification \hfill\normalfont\footnotesize E2 disc. 2021 (Q13), Rapport jury 2021}
]
\textbf{Énoncé.}
5 onduleurs à 6 entrées MPPT = 30 groupes. Répartir 782 panneaux. Calculer $V_{MPP}$, $I_{MPP}$ et $P_{MPP}$ par entrée. Vérifier la compatibilité avec l'onduleur.

\tcbline

\textbf{Corrigé.}

\textbf{Étape 1 — Panneaux par entrée.}
\[
  N_{s/groupe} = \frac{782}{30} \approx 26\,\text{panneaux par entrée}
\]

\textbf{Étape 2 — Grandeurs par entrée.}
\[
  V_{MPP,entree} = 26 \times 34{,}32 = \mathbf{892\,\mathrm{V}}
  \qquad;\qquad
  I_{MPP,entree} = 9{,}34\,\mathrm{A}
  \qquad;\qquad
  P_{MPP,entree} = 26 \times 320 = \mathbf{8320\,\mathrm{W}}
\]

\textbf{Conclusion.} Les valeurs calculées ($V_{MPP} = 892\,\mathrm{V}$, $I_{MPP} = 9{,}34\,\mathrm{A}$) doivent être comparées aux plages min-max de l'onduleur sur sa datasheet constructeur avant de valider le dimensionnement : le jury insiste sur cette vérification explicite, pas seulement sur le calcul.
\end{tcolorbox}

\vspace{4pt}

% -- QT3 -------------------------------------------------------------

\begin{tcolorbox}[questiontype,
  title={QT3 - Algorithme MPPT - perturbation et observation \hfill\normalfont\footnotesize E2 disc. 2025 (Q17-Q18), Rapport jury 2025}
]
\textbf{Énoncé.}
Sur la courbe $P = f(U)$ d'un champ PV, le point MPP est à $U_{MPP} = 317\,\mathrm{V}$.\\
\textbf{Cas P1 :} $U_1 = 195\,\mathrm{V}$, on impose $\Delta U > 0$. Quel est le signe de $\Delta P$ ? Que fait le régulateur MPPT ?\\
\textbf{Cas P2 :} $U_2 = 340\,\mathrm{V}$, on impose $\Delta U > 0$. Même question.

\tcbline

\textbf{Corrigé.}

\textbf{Cas P1 ($U_1 = 195\,\mathrm{V} < U_{MPP}$).} On est à gauche du MPP. En augmentant $U$, on se rapproche du MPP - $\mathbf{\Delta P > 0}$.
Réponse du régulateur : continuer dans le même sens (ajuster le rapport cyclique du hacheur pour continuer d'augmenter $U$).

\textbf{Cas P2 ($U_2 = 340\,\mathrm{V} > U_{MPP}$).} On est à droite du MPP. En augmentant $U$, on s'éloigne - $\mathbf{\Delta P < 0}$.
Réponse du régulateur : inverser le sens (ajuster le rapport cyclique en sens inverse pour diminuer $U$).

\textbf{Conclusion.} Le processus converge quand $\Delta P \approx 0$, c'est-à-dire au voisinage de $U_{MPP} = 317\,\mathrm{V}$ : c'est ce critère d'arrêt, et non une valeur de tension fixée à l'avance, qui pilote l'algorithme P\&O en continu.
\end{tcolorbox}

\end{tcolorbox}

% -----------------------------------------------------------------------
%  4. ATTENTES DU JURY
% -----------------------------------------------------------------------

\begin{tcolorbox}[
  jury,
  title={4. Attentes du jury - Erreurs récurrentes et points de vigilance},
]

\begin{enumerate}[label=\textbf{\arabic*.}, leftmargin=*, itemsep=5pt]

  \item \textbf{Correction en température obligatoire.}
  À $50^\circ\mathrm{C}$, multiplier la puissance STC par $1 - 0{,}004 \times (50 - 25) = 0{,}90$. Rapport jury 2021 : «~Le jury déplore que certains candidats se limitent à donner des résultats sans réelle explication et à ne pas répondre entièrement aux questions lorsqu'il est demandé de conclure sur le choix d'un composant.~»

  \item \textbf{Vérifier la compatibilité avec l'onduleur.}
  Après avoir calculé $V_{MPP}$ et $I_{MPP}$, vérifier explicitement que ces valeurs sont dans les plages admissibles de l'onduleur. Omission signalée comme lacune fréquente.

  \item \textbf{Cellules en série dans un panneau $\neq$ panneaux en série.}
  Pour une cellule dans un panneau de 60 cellules série : $I_{sc,cell} = I_{sc,panneau}$ mais $V_{oc,cell} = V_{oc,panneau}/60$. Ne pas confondre avec l'association de panneaux.

  \item \textbf{MPPT : sens de la perturbation.}
  Si $\Delta P > 0$ après $+\Delta U$ - continuer dans le même sens. Si $\Delta P < 0$ - inverser. Beaucoup de candidats inversent la logique.

  \item \textbf{Puissance éolienne $\propto v^3$.}
  Si la vitesse du vent double, la puissance est multipliée par $2^3 = 8$, pas par 2. Erreur fréquente.

  \item \textbf{Rendement STEP : conclure sur l'intérêt contextuel.}
  Citer le chiffre numérique (58\,\%) ET la conclusion : "peu performant, mais écologique, préférable à des batteries supplémentaires, exploitable géographiquement".

\end{enumerate}

\medskip
\begin{tcolorbox}[
  colback=colJuryFrame!12!colJury,
  colframe=colJuryFrame,
  left=4pt, right=4pt, top=2pt, bottom=2pt,
  fontupper=\small\itshape
]
Rapport jury 2021 : «~Même si quasiment tous les candidats (98\,\%) ont traité cette partie, la plupart d'entre eux ont rencontré des difficultés à justifier le choix de l'onduleur.~»\\
Rapport jury 2025 : «~MPPT : cette commande permet d'extraire le maximum de puissance. Le rapport cyclique augmente ou diminue en fonction de l'évolution de dP jusqu'à obtenir le point maximum (dP = 0).~»
\end{tcolorbox}

\end{tcolorbox}

% -----------------------------------------------------------------------
%  PIED - PRESENCE DANS LES SUJETS
% -----------------------------------------------------------------------

\begin{tcolorbox}[
  ficheSection=colPresence,
  title={Présence dans les sujets analysés (2013-2026)},
]
\begin{center}
\footnotesize\renewcommand{\arraystretch}{1.15}
\begin{tabular}{lccccccccccccc}
\toprule
\textbf{Année} & 13 & 14 & 15 & 16 & 17 & 18 & 19 & 20 & 21 & 22 & 23 & 24 & 25 \\
\midrule
E1 (commune)   & - & - & - & - & - & - & - & - & - & - & - & - & - \\
E2 discipl. IE & - & - & - & $\checkmark$ & - & - & - & $\checkmark$ & $\checkmark$ & $\checkmark$ & $\checkmark$ & - & $\checkmark$ \\
\bottomrule
\end{tabular}
\end{center}
\smallskip
\footnotesize\textbf{Sous-thèmes fréquents :}
panneau PV - nombre et superficie (4/13) - MPPT algorithme P\&O (3/13) - dimensionnement onduleur (3/13) - micro-hydraulique énergie pot. (2/13) - comparaison filières renouvelables (2/13)

\medskip
\textbf{Ordres de grandeur (repères rapides).}
\begin{itemize}[leftmargin=1.3em]
  \item Réseau BT FR : tension phase-neutre typique \SI{230}{\volt}, tension phase-phase typique \SI{400}{\volt}, fréquence \SI{50}{\hertz}.
  \item Courants d'étude fréquents en sujet : de quelques \si{\ampere} à quelques dizaines de \si{\ampere}.
  \item En conversion d'énergie, un rendement attendu reste généralement entre 0{,}8 et 0{,}98 selon la technologie.
\end{itemize}

\medskip
\textbf{Synthèse express (à refaire sans support).}
\begin{enumerate}[leftmargin=1.4em]
  \item Poser les hypothèses du modèle et les conventions de signe.
  \item Écrire 2 à 4 relations clés avec unités.
  \item Réaliser un calcul numérique complet et vérifier l'ordre de grandeur.
  \item Justifier le résultat par un argument technique de niveau CAPET.
\end{enumerate}

\end{tcolorbox}

\CAPETENDFILE
